\chapter{US Equity Market Structure}\label{ch:m1:us-equity-market-structure}

One hundred shares of a large American company can be bought, at this
instant, on more than a dozen stock exchanges, in dozens of \omterm{def:m1:us-equity-market-structure:dark}{dark pools} and
from dealers who never show a quote to anyone. More than half of
the 17.6 billion shares that changed hands on an average day of 2025 did so
on none of the exchanges. For two decades a single regulation has tied these
places together by decreeing which of their prices is ``the'' price and
forbidding trades at worse ones; in June 2026 the regulator proposed to
repeal that decree. This chapter describes the machine as it stands: the
rules, the consolidated quote they create, the two speeds at which that quote
can be known, the fee games played around it, and the half of the market that
lives off-exchange.

\section{Regulation NMS}

The Securities and Exchange Commission adopted Regulation National Market
System in 2005. Four of its rules shape every order.

\begin{definition}[Protected quotation, trade-through, order protection rule]\label{def:m1:us-equity-market-structure:opr}
A \emph{protected quotation} is the best bid or best offer of an exchange
that is displayed, automated and immediately accessible. A
\emph{trade-through}\index{trade-through} is the execution of a trade at a
price inferior to a protected quotation displayed elsewhere: a purchase above
another venue's protected offer, or a sale below its protected bid. The
\emph{order protection rule}\index{order protection rule} (Rule 611) requires
every trading centre to have policies reasonably designed to prevent
trade-throughs.
\end{definition}

\begin{definition}[Access fee cap]\label{def:m1:us-equity-market-structure:cap}
The \emph{access fee cap}\index{access fee cap} (Rule 610(c)) limits what a
venue may charge for executing against its protected quotation: \$0.003 a
share (``30 mils'') for stocks priced at a dollar or more. Without it a venue
could display an attractive price that the protection rule forces others to
route to, and charge whatever it liked for reaching it.
\end{definition}

Rule 610 also forbids quotations that \emph{lock} or \emph{cross} the market
(a bid equal to, or above, another venue's offer). Rule 612, the
\emph{sub-penny rule}, fixes the minimum quoting increment at one cent for
stocks at or above a dollar. Rule 603 governs the distribution of market
data, to which we return below.

\begin{dated}{2026-09}{Regulation NMS in motion}\label{dat:m1:us-equity-market-structure:nms}
In September 2024 the SEC adopted amendments creating a half-cent quoting
increment for ``tick-constrained'' stocks and cutting the \omterm{def:m1:us-equity-market-structure:cap}{access fee cap} to
\$0.001. Neither is in force: by an order of 11 June 2026 compliance with
both is deferred to the first business day of November 2027. On the same day
the SEC \emph{proposed to rescind} Rule 611 and the ban on locked and crossed
markets altogether, with a 60-day comment period. The rest of this chapter
describes the rules in force; read the proposal's current status before
building anything that depends on them.
\end{dated}

\section{The national best bid and offer}

\begin{definition}[Round lot and odd lot]\label{def:m1:us-equity-market-structure:lots}
A \emph{round lot}\index{round lot} is the standard trading unit that a
quotation must reach to be protected; an \emph{odd lot}\index{odd lot} is any
smaller quantity.
\end{definition}

\begin{dated}{2026-09}{Round lots by price}\label{dat:m1:us-equity-market-structure:lots}
Since 3 November 2025 the \omterm{def:m1:us-equity-market-structure:lots}{round lot} depends on the stock's average closing
price over the previous evaluation period: 100 shares up to \$250.00; 40
shares from \$250.01 to \$1\,000.00; 10 shares from \$1\,000.01 to
\$10\,000.00; 1 share above. From the same date the consolidated tapes quote
sizes in shares instead of lots. Until then a bid for 50 shares of a
\$3\,000 stock, worth \$150\,000, was an \omterm{def:m1:us-equity-market-structure:lots}{odd lot}: neither protected nor
included in the NBBO.
\end{dated}

\begin{definition}[National best bid and offer]\label{def:m1:us-equity-market-structure:nbbo}
The \emph{national best bid and offer}\index{national best bid and offer}
(NBBO) of a stock is the highest protected bid and the lowest protected
offer across all exchanges, with the total size displayed at each of those
prices.
\end{definition}

\begin{example}[Three venues]\label{ex:m1:us-equity-market-structure:nbbo}
Venue X shows $10.00 \times 10.02$ (300 by 500 shares), venue Y $10.01
\times 10.02$ (200 by 100), venue Z $10.01 \times 10.03$ (400 by 900). The
NBBO is $10.01 \times 10.02$, 600 by 600: the bid from Y and Z, the offer
from X and Y. A broker that fills a client's market purchase at 10.03 while
X and Y offer at 10.02 has traded through two protected quotations. If a
fourth venue bids 10.015 for 60 shares, nothing changes: an \omterm{def:m1:us-equity-market-structure:lots}{odd lot} is not
protected. The NBBO is the benchmark of every execution-quality statistic of
\cref{ch:m1:retail-flow-and-wholesaling}; it is also, as we now see, not a
single thing.
\end{example}

\section{One market, two clocks}

\begin{definition}[Securities information processor and direct feed]\label{def:m1:us-equity-market-structure:sip}
A \emph{securities information processor}\index{securities information processor}
(SIP) collects the best quotes and the trades of every exchange, computes the
NBBO, and publishes the consolidated tape. A
\emph{direct feed}\index{direct feed} is the proprietary data stream sold by
an exchange, carrying its full order book, from which a subscriber can
compute its own consolidated view.
\end{definition}

\begin{omfigure}
\begin{tikzpicture}[font=\small,
  ven/.style={draw=omExo,rounded corners=2pt,minimum width=2.1cm,minimum height=0.7cm,fill=omExo!6},
  box/.style={draw=omDef,rounded corners=2pt,minimum width=3.2cm,minimum height=0.9cm,align=center,fill=omDef!4},
  slow/.style={-{Stealth[length=2mm]},thick,omThm},
  fast/.style={-{Stealth[length=2mm]},thick,omDef}]
\foreach \i/\y/\n in {a/1.8/Exchange A,b/0.6/Exchange B,c/-0.6/Exchange C,d/-1.8/Exchange D}{ \node[ven] (v\i) at (0,\y) {\n}; }
\node[box] (sip) at (5.6,1.5) {SIP\\(consolidates, republishes)};
\node[box] (own) at (5.6,-1.5) {Firm's own book builder\\in the same data centres};
\node[box] (slowu) at (11.6,1.5) {Tape subscriber\\sees the NBBO};
\node[box] (fastu) at (11.6,-1.5) {Direct-feed subscriber\\sees it earlier};
\foreach \i in {a,b,c,d}{ \draw[slow] (v\i.east) -- (sip.west); \draw[fast] (v\i.east) -- (own.west); }
\draw[slow] (sip) -- node[above,font=\footnotesize]{$+\,\delta$} (slowu);
\draw[fast] (own) -- (fastu);
\end{tikzpicture}

\omcaption{The same quotes, consolidated twice. The tape's path (red) adds a
delay $\delta$: the trip to the processor, its computation, and the trip to
the subscriber. The direct path (blue) is as fast as the subscriber's own
lines and code.}\label{fig:m1:us-equity-market-structure:twoclocks}
\end{omfigure}

\begin{proposition}[How often the tape is wrong]\label{prop:m1:us-equity-market-structure:stale}
Suppose the best bid changes at the times of a Poisson process of rate
$\lambda$, and the tape reports each change $\delta$ later than a direct-feed
reader observes it. If $\lambda\delta \ll 1$, the fraction of time during
which the two views differ is approximately $\lambda\delta$.
\end{proposition}
\begin{proof}
Each change opens an interval of length $\delta$ during which the tape still
shows the old bid. When $\lambda\delta \ll 1$ these intervals overlap with
negligible probability, and over a period $T$ they cover $\lambda T \cdot
\delta$ in expectation.
\end{proof}

\begin{example}[Half a percent of the time, all of the trades that matter]\label{ex:m1:us-equity-market-structure:halfpercent}
With five changes a second and $\delta = 1$ millisecond, the tape is stale
0.5\% of the time; in a burst of fifty changes a second, 5\%
(\cref{fig:m1:us-equity-market-structure:stale}). The fraction looks
negligible and is not: these are precisely the moments when the price is
moving, and an order priced off the stale view is filled \emph{because} it is
stale. A quote that has not yet followed a two-cent move can be bought one
cent below the new fair value; in the tutorial's simulated stock the slower
venues offer 3\,636 such opportunities in ten minutes. Who is allowed to take
them, and how venues defend their slow participants, is the business of One
Quant Book 11.
\end{example}

\begin{omfigure}
\begin{tikzpicture}
\begin{axis}[width=10.5cm,height=4.8cm,
  xlabel={delay of the tape $\delta$ (microseconds)},ylabel={time with a stale best bid (\%)},
  tick label style={font=\small},label style={font=\small},
  xmin=0,xmax=2000,ymin=0,ymax=10,grid=major,grid style={gray!25},xtick={0,500,1000,1500,2000},
  xticklabel style={/pgf/number format/1000 sep={\,}},
  legend columns=2,legend style={font=\footnotesize,at={(0.5,-0.34)},anchor=north,draw=none,fill=none,/tikz/every even column/.append style={column sep=8pt}},
  table/col sep=comma]
\addplot[omThm,very thick,mark=*,mark size=1.2pt] table[x=sip_delay_us,y=stale_pct_busy]{figdata/markets-1/09-us-equity-market-structure/stale.csv};
\addplot[omDef,very thick,mark=*,mark size=1.2pt] table[x=sip_delay_us,y=stale_pct_quiet]{figdata/markets-1/09-us-equity-market-structure/stale.csv};
\legend{50 price changes a second,5 price changes a second}
\end{axis}
\end{tikzpicture}

\omcaption{Staleness of a consolidated best bid against the tape's delay, for
four venues repricing 50 to 400 microseconds after each move. The slopes are
the two values of $\lambda$. Data: the tutorial's
simulation.}\label{fig:m1:us-equity-market-structure:stale}
\end{omfigure}

\section{Exchange families and the fee game}

Most exchanges belong to a few groups, each operating several venues with
different fee schedules on nearly identical technology. The reason is
the \omterm{def:m1:us-equity-market-structure:cap}{access fee cap}: venues cannot compete on the displayed price, which the
tick fixes at a cent, so they compete inside the cap, on who pays whom.

\begin{definition}[Maker-taker pricing and inverted venue]\label{def:m1:us-equity-market-structure:makertaker}
Under \emph{maker-taker pricing}\index{maker-taker pricing} a venue charges
the order that removes liquidity (the \emph{taker}) a fee and pays the
resting order it executed against (the \emph{maker}) a smaller rebate,
keeping the difference. An \emph{inverted venue}\index{inverted venue} does
the opposite: it pays the taker and charges the maker.
\end{definition}

\begin{proposition}[Fees are a finer tick]\label{prop:m1:us-equity-market-structure:finer}
Let two venues display the same offer $a$. A buyer taking liquidity pays
$a + \varphi$ on a venue with taker fee $\varphi$ and $a - \rho$ on an
\omterm{def:m1:us-equity-market-structure:makertaker}{inverted venue} with taker rebate $\rho$. A rational taker therefore hits the
\omterm{def:m1:us-equity-market-structure:makertaker}{inverted venue} first: a \omterm{def:m1:us-equity-market-structure:makertaker}{maker} who posts there is at the front of a
\emph{cross-venue} queue, for a price, and receives $a - \mu$ with $\mu$ the
\omterm{def:m1:us-equity-market-structure:makertaker}{inverted venue}'s \omterm{def:m1:us-equity-market-structure:makertaker}{maker} fee.
\end{proposition}
\begin{proof}
Immediate from the definitions; the ordering of venues by all-in price
$a + \text{fee}$ is strict although the displayed prices are equal.
\end{proof}

\begin{example}[Buying priority]\label{ex:m1:us-equity-market-structure:priority}
With illustrative fees inside the cap: on a maker-taker venue a seller
resting at 10.00 earns a rebate of 0.20 cent but waits behind everyone who
arrived earlier; on an \omterm{def:m1:us-equity-market-structure:makertaker}{inverted venue} it pays 0.18 cent and is executed ahead
of every maker-taker venue, because takers collect 0.15 cent for going there
first. The difference, 0.38 cent, is the price of jumping the queue in a
market whose tick is one cent. A \omterm{def:m1:what-a-trading-firm-does:market-maker}{market maker} chooses per stock and per
moment: where the queue is long and the spread is one tick, priority is
worth buying (\cref{fig:m1:us-equity-market-structure:fees}).
\end{example}

\begin{omfigure}
\begin{tikzpicture}
\begin{axis}[width=9.2cm,height=3.8cm,xbar,bar width=8pt,
  xlabel={all-in cost of taking an offer, relative to the quote (cents)},
  ytick=data,yticklabels from table={figdata/markets-1/09-us-equity-market-structure/fees.csv}{label},
  yticklabel style={font=\small},tick label style={font=\small},label style={font=\small},
  y dir=reverse,xmin=-0.35,xmax=0.45,xtick={-0.3,-0.15,0,0.15,0.3,0.45},enlarge y limits=0.25,grid=major,grid style={gray!25},
  xticklabel style={/pgf/number format/fixed,/pgf/number format/precision=2},
  nodes near coords,nodes near coords style={font=\footnotesize,/pgf/number format/fixed,/pgf/number format/precision=2},
  table/col sep=comma]
\addplot[fill=omDef!60,draw=omDef] table[x=net_cents_vs_quote,y=k]{figdata/markets-1/09-us-equity-market-structure/fees.csv};
\end{axis}
\end{tikzpicture}

\omcaption{Three venues displaying the same offer are three different prices
to a taker. Fee levels are illustrative, within the 0.30-cent cap. Data: the
chapter's script.}\label{fig:m1:us-equity-market-structure:fees}
\end{omfigure}

\begin{definition}[Intermarket sweep order]\label{def:m1:us-equity-market-structure:iso}
An \emph{intermarket sweep order}\index{intermarket sweep order} (ISO) is a
limit order marked so that the receiving venue may execute it at once without
checking other venues' protected quotations, because the sender has
simultaneously routed orders to take out every better-priced protected
quotation elsewhere. Compliance with the \omterm{def:m1:us-equity-market-structure:opr}{order protection rule} moves from the
venue to the sender.
\end{definition}

The ISO is how every large or fast participant actually trades: it sends, in
the same microsecond, one order per venue, sized to the displayed quantity,
and so sweeps several price levels across the market before any venue can
re-route anything. Without it an order for more than the size at the best
price would have to crawl through the venues one routing decision at a time.

\section{The off-exchange half}

\begin{definition}[Dark pool]\label{def:m1:us-equity-market-structure:dark}
A \emph{dark pool}\index{dark pool} is an \omterm{def:m1:exchanges-brokers-venues:ats}{alternative trading system}
(\cref{def:m1:exchanges-brokers-venues:ats}) that displays no quotations:
orders rest unseen and execute against each other at prices derived from the
NBBO, most often its midpoint.
\end{definition}

Trades that occur off-exchange are still public: they are reported within
seconds to a \emph{trade reporting facility} and appear on the tape without
the name of the place. They come from two sources. \omterm{def:m1:us-equity-market-structure:dark}{Dark pools} serve
institutions that do not want to reveal a large order and are content to
wait for a match at the midpoint. \emph{Principal dealers} --- wholesalers
executing retail brokers' orders against their own inventory
(\cref{ch:m1:retail-flow-and-wholesaling}) and banks' single-dealer platforms
--- trade as counterparty, not as venue.

\begin{dated}{2026-09}{Where US shares traded in 2025}\label{dat:m1:us-equity-market-structure:where}
Average daily volume in 2025 was 17.6 billion shares worth \$1.1 trillion, up
about 45\% on 2024. The trade reporting facilities' share of consolidated
volume reached 50.6\%, above one half for the first time over a full year;
81.3\% of that off-exchange volume was executed by principal dealers and
18.7\% on \omterm{def:m1:exchanges-brokers-venues:ats}{alternative trading systems}, according to one exchange group's
annual review.
\end{dated}

\begin{omfigure}
\begin{tikzpicture}
\begin{axis}[width=9cm,height=3.8cm,xbar,bar width=8pt,
  xlabel={share of 2025 US equity volume (\%)},
  ytick=data,yticklabels from table={figdata/markets-1/09-us-equity-market-structure/where_2025.csv}{label},
  yticklabel style={font=\small},tick label style={font=\small},label style={font=\small},
  y dir=reverse,xmin=0,xmax=60,enlarge y limits=0.25,grid=major,grid style={gray!25},
  nodes near coords,nodes near coords style={font=\footnotesize},
  table/col sep=comma]
\addplot[fill=omDef!60,draw=omDef] table[x=pct,y=k]{figdata/markets-1/09-us-equity-market-structure/where_2025.csv};
\end{axis}
\end{tikzpicture}

\omcaption{Where a share traded in 2025. The exchanges, which produce the
quotes that price everything else, executed a minority of the volume. Data:
as in \cref{dat:m1:us-equity-market-structure:where}.}
\end{omfigure}

\begin{remark}[Why the split worries regulators]\label{rem:m1:us-equity-market-structure:free}
Off-exchange trades are priced from the NBBO and contribute nothing to it:
they free-ride on displayed quotes. If the orders least dangerous to a \omterm{def:m1:what-a-trading-firm-does:market-maker}{market
maker} --- retail orders --- are executed away from exchanges, the orders
that do reach exchanges are on average better informed, exchange spreads
widen (\cref{prop:m1:what-a-trading-firm-does:capture}), and the benchmark
against which off-exchange ``price improvement'' is measured becomes easier
to beat. Whether this loop is large in practice is one of the most contested
empirical questions in market structure.
\end{remark}

\section{Speed bumps}

If \cref{prop:m1:us-equity-market-structure:stale} describes a harm, a venue
can sell protection from it. In 2016 the SEC approved as an exchange a venue,
IEX, whose every inbound and outbound message passes through 38 miles of
coiled optical fibre, a delay of 350 microseconds, while the venue's own view
of the market is \emph{not} delayed. Resting orders pegged to the midpoint
are therefore repriced by the venue before a fast trader who has seen the
same move can reach them. The Commission held that a delay of this size is
\emph{de minimis}: the venue's quotations remain ``immediately accessible''
and protected. The debate about who should bear the cost of speed is far from
closed.

\section{Tutorial: building an NBBO and timing the tape}\label{tut:m1:us-equity-market-structure:nbbo}

\begin{tutorial}
\textbf{Goal.} Consolidate per-venue quotes into an NBBO, flag locked and
crossed markets and \omterm{def:m1:us-equity-market-structure:opr}{trade-throughs}, then measure how stale a delayed tape is.
\textbf{End state:} the numbers of
\cref{ex:m1:us-equity-market-structure:nbbo} and
\cref{fig:m1:us-equity-market-structure:stale}.
\begin{tutsteps}
\item \textbf{The builder}, as it runs inside the firm's trading process: a
fixed table of venues, no allocation, one pass per update.
\omcode{code/firm/nbbo/cpp/firm_nbbo.hpp}{41}{58}{Recomputing the NBBO from the venue table: protected sizes only, ties pooled.}\label{lst:m1:us-equity-market-structure:cpp}
\item \textbf{Check it on three venues.} The acceptance tests replay
\cref{ex:m1:us-equity-market-structure:nbbo}: best bid 10.01 for 600 shares
on two venues; an odd-lot bid of 10.015 ignored until it reaches a \omterm{def:m1:us-equity-market-structure:lots}{round lot};
a purchase at 10.03 flagged as a \omterm{def:m1:us-equity-market-structure:opr}{trade-through}.
\item \textbf{Simulate two clocks.} Jumps of one to three ticks arrive at
Poisson times; each venue reprices after its own delay; the tape adds
$\delta$.
\omcode{code/markets-1/09-us-equity-market-structure/python/stale_tape.py}{29}{49}{Each venue's quote as a step function of time, and the measure of the set where tape and direct view differ.}\label{lst:m1:us-equity-market-structure:stale}
\item \textbf{Compare with the proposition.} At $\delta = 1\,000$
microseconds and $\lambda = 5$ per second the simulation gives 0.51\%
against a predicted 0.50\%; at $\lambda = 50$, 4.9\% against 5\%, the
shortfall being the overlaps the proposition neglects.
\end{tutsteps}
\textbf{What to change next.} Give one venue a 350-microsecond speed bump on
incoming orders only, and count how many of its stale quotes survive. Then
add an \omterm{def:m1:us-equity-market-structure:makertaker}{inverted venue} and route a taker by all-in price.
\end{tutorial}

\section{Build: the consolidated quote}\label{bld:m1:us-equity-market-structure:nbbo}

\begin{build}
\bfield{Purpose} Every strategy, router and execution report of the
miniature firm needs the best bid and offer across venues, computed by the
firm itself from \omterm{def:m1:us-equity-market-structure:sip}{direct feeds} (the feed handlers are built in One Quant Book
13).
\bfield{Interface} \texttt{NbboBuilder(round\_lot)};
\texttt{update(venue, bid, bid\_size, ask, ask\_size)} returning the new NBBO
if it changed; \texttt{nbbo()}; \texttt{trades\_through(side, price)}. The
NBBO carries prices, pooled sizes, the set of venues at each price, and
\texttt{locked} and \texttt{crossed} flags.
\bfield{Rules} Integer prices. A quote smaller than the \omterm{def:m1:us-equity-market-structure:lots}{round lot} is stored
but not protected. A missing side means no NBBO. In C++ and Rust: a fixed
venue table, no allocation and no locking on the update path; venues are
small integers and the venue sets are bit masks.
\bfield{Acceptance tests} \texttt{code/firm/nbbo/tests/}, \texttt{cpp/} and
\texttt{rust/}: the three-venue example; \omterm{def:m1:us-equity-market-structure:lots}{odd lots}; change detection; locked
and crossed; \omterm{def:m1:us-equity-market-structure:opr}{trade-through}; rejection of malformed quotes.
\bfield{Stretch} Per-symbol \omterm{def:m1:us-equity-market-structure:lots}{round lots} from a reference table
(\cref{dat:m1:us-equity-market-structure:lots}); a second, \emph{fee-adjusted}
best bid and offer using each venue's taker fee
(\cref{prop:m1:us-equity-market-structure:finer}).
\end{build}

\begin{omsources}
\item US Securities and Exchange Commission, \emph{Regulation NMS}, Release
34-51808, \emph{Federal Register} 70 (29 June 2005).
\item US Securities and Exchange Commission, \emph{Fact Sheet: Regulation NMS
Reforms} (proposal of 11 June 2026, Release 34-105655); Release 34-105656,
order granting temporary exemptive relief, 11 June 2026; press release
2024-137 (amendments of 18 September 2024).
\item Cboe Global Markets, \emph{Regulation NMS Round Lots Enhancements FAQ},
28 October 2025; \emph{2025 U.S. Equities Year in Review}.
\item US Securities and Exchange Commission, Release 34-78102,
\emph{Commission Interpretation Regarding Automated Quotations Under
Regulation NMS}, June 2016; IEX, \emph{The speed bump}.
\item E.~Budish, P.~Cramton and J.~Shim, ``The high-frequency trading arms
race'', \emph{Quarterly Journal of Economics} 130 (2015).
\item M.~O'Hara, ``High frequency market microstructure'', \emph{Journal of
Financial Economics} 116 (2015).
\end{omsources}

\section{Exercises}

\begin{exercise}[$\star$]\label{exo:m1:us-equity-market-structure:1}
Four venues quote a stock with a 100-share \omterm{def:m1:us-equity-market-structure:lots}{round lot}: A $25.10 \times 25.14$
(500 by 200), B $25.11 \times 25.13$ (100 by 300), C $25.11 \times 25.13$ (300
by 50), D $25.12 \times 25.15$ (80 by 400). Give the NBBO with sizes and
venues.
\end{exercise}

\begin{exercise}[$\star$]\label{exo:m1:us-equity-market-structure:2}
Using \cref{dat:m1:us-equity-market-structure:lots}, give the \omterm{def:m1:us-equity-market-structure:lots}{round lot} of
stocks priced at \$48, \$251, \$999, \$4\,200 and \$712\,000, and the dollar
value of one \omterm{def:m1:us-equity-market-structure:lots}{round lot} of each.
\end{exercise}

\begin{exercise}[$\star$]\label{exo:m1:us-equity-market-structure:3}
From \cref{dat:m1:us-equity-market-structure:where}, compute the average
daily number of shares executed on exchanges, by principal dealers and on
\omterm{def:m1:exchanges-brokers-venues:ats}{alternative trading systems}, and the average price of a share traded.
\end{exercise}

\begin{exercise}[$\star\star$]\label{exo:m1:us-equity-market-structure:4}
With the NBBO of exercise 1, a broker executes a client's purchase of 200
shares at 25.14 on venue A. Which protected quotations were traded through?
What should the broker have sent, and marked how, to buy 500 shares
immediately and lawfully?
\end{exercise}

\begin{exercise}[$\star\star$]\label{exo:m1:us-equity-market-structure:5}
A maker-taker venue charges takers 0.30 cent and pays \omterm{def:m1:us-equity-market-structure:makertaker}{makers} 0.25; an
\omterm{def:m1:us-equity-market-structure:makertaker}{inverted venue} pays takers 0.10 and charges \omterm{def:m1:us-equity-market-structure:makertaker}{makers} 0.14. Both show an offer
of 40.00. Give the all-in price paid by a taker and received by a \omterm{def:m1:us-equity-market-structure:makertaker}{maker} on
each, the venue's revenue per share on each, and the price of priority.
\end{exercise}

\begin{exercise}[$\star\star$]\label{exo:m1:us-equity-market-structure:6}
A stock's best bid changes 12 times a second on average. Estimate the
fraction of time a tape delayed by 600 microseconds shows a stale bid, and
the number of seconds per 6.5-hour trading day. Why does the proposition
overestimate slightly when changes cluster?
\end{exercise}

\begin{exercise}[$\star\star\star$]\label{exo:m1:us-equity-market-structure:7}
\emph{Coding.} With \texttt{stale\_tape}, halve every venue's repricing delay
and recompute the stale fraction at $\delta = 500$ microseconds and $\lambda =
5$. Explain why the answer barely changes, and what quantity the venues'
delays \emph{do} control.
\end{exercise}

\begin{exercise}[$\star\star\star$]\label{exo:m1:us-equity-market-structure:8}
\emph{Find the flaw.} A retail broker's report states: ``98\% of our clients'
orders were executed at or better than the NBBO.'' The broker computes the
NBBO from the consolidated tape and timestamps each order on receipt at its
own servers. Give two reasons why the statistic overstates execution quality,
and one reason specific to high-priced stocks before November 2025.
\end{exercise}

\section{Problem: A Stale Tape}

\begin{problem}[{Weekend problem --- two views of one stock}]\label{pb:m1:us-equity-market-structure:1}
A stock trades on four exchanges. Each quotes one tick (one cent) either side
of its estimate of the fair price, for 300 shares a side. When the fair price
jumps, the exchanges' \omterm{def:m1:what-a-trading-firm-does:market-maker}{market makers} reprice after 60, 150, 250 and 500
microseconds respectively. The consolidated tape reaches a subscriber 900
microseconds after a \omterm{def:m1:us-equity-market-structure:sip}{direct feed} would. The fair price jumps 8 times a minute
on average; 65\% of jumps are of one tick, 25\% of two ticks and 10\% of three.

\medskip\noindent\textbf{Part I --- The tape.}
\begin{enumerate}
\item For what fraction of the time does the tape show a stale best bid?
\item How many seconds is that in a 6.5-hour day?
\item After an upward jump, which venue's update moves the best bid on a
\omterm{def:m1:us-equity-market-structure:sip}{direct feed}, and after how long?
\item After the same jump, when does the best \emph{offer} move, and why is
the answer different?
\end{enumerate}

\medskip\noindent\textbf{Part II --- The snipe.}
A fast trader sees every \omterm{def:m1:us-equity-market-structure:sip}{direct feed} and can reach any venue 40 microseconds
after a jump.
\begin{enumerate}[resume]
\item After an upward jump of $J$ ticks, which venues still show the old
offer when the trader arrives?
\item Buying at an old offer and valuing at the new fair price, what is the
profit per share for $J = 1, 2, 3$?
\item Give the expected profit per share per jump, per stale venue.
\item With 300 shares on each stale venue, give the expected profit per
jump and per day if the trader wins every race.
\item The trader pays a taker fee of 0.30 cent. Recompute the expected profit
per share per jump per venue. Which jumps are still worth taking?
\end{enumerate}

\medskip\noindent\textbf{Part III --- The defence.}
\begin{enumerate}[resume]
\item The slowest \omterm{def:m1:what-a-trading-firm-does:market-maker}{market maker} loses what the trader wins. What does being
sniped cost it per day, and per share it trades if it trades 400\,000 shares
a day?
\item Its \omterm{prop:m1:what-a-trading-firm-does:capture}{net capture} before sniping was 0.35 cent a share. Is it still
profitable?
\item It can cut its delay to 30 microseconds for \$2.5 million a year. Is
the investment worth it, on 252 days?
\item If all four \omterm{def:m1:what-a-trading-firm-does:market-maker}{market makers} make that investment, what happens to the
trader's profit and to the \omterm{def:m1:what-a-trading-firm-does:market-maker}{market makers}' costs? What is this situation
called?
\item Alternatively the venue imposes a 350-microsecond delay on incoming
orders. Which of its \omterm{def:m1:what-a-trading-firm-does:market-maker}{market makers}' quotes are now safe?
\end{enumerate}

\medskip\noindent\textbf{Part IV --- The rule.}
\begin{enumerate}[resume]
\item A broker routes a client's buy order using the tape. During a stale
interval after an upward jump, where does it send the order, and what
happens?
\item Is the fill a \omterm{def:m1:us-equity-market-structure:opr}{trade-through} under Rule 611? Explain with reference to
what each venue displays at the time.
\item The SEC's 2024 amendments would cut the \omterm{def:m1:us-equity-market-structure:cap}{access fee cap} to 0.10 cent.
Redo question 9.
\item Would rescinding Rule 611 remove the sniping opportunity?
\item State the \emph{named result}: the trader's expected profit per share
from a stale quote, per jump, net of the current taker fee, in cents.
\item In one sentence, who pays for it in the end?
\end{enumerate}
\end{problem}

\section{Interview questions}

\begin{interviewq}[$\star$ \iqroles{trader, developer, researcher}]\label{iq:m1:us-equity-market-structure:1}
What is the NBBO, who computes it, and why might two firms disagree about
what it was at a given microsecond?
\end{interviewq}

\begin{interviewq}[$\star$ \iqroles{trader, developer}]\label{iq:m1:us-equity-market-structure:2}
What is an \omterm{def:m1:us-equity-market-structure:iso}{intermarket sweep order} and why does it exist?
\end{interviewq}

\begin{interviewq}[$\star\star$ \iqroles{trader, researcher}]\label{iq:m1:us-equity-market-structure:3}
Why would anyone post liquidity on a venue that charges \omterm{def:m1:us-equity-market-structure:makertaker}{makers} a fee when
another venue pays them a rebate?
\end{interviewq}

\begin{interviewq}[$\star\star$ \iqroles{developer}]\label{iq:m1:us-equity-market-structure:4}
You must maintain the NBBO for 8\,000 symbols from sixteen \omterm{def:m1:us-equity-market-structure:sip}{direct feeds} with
minimal latency. Describe your data structures and threading model.
\end{interviewq}

\begin{interviewq}[$\star\star$ \iqroles{researcher, trader}]\label{iq:m1:us-equity-market-structure:5}
More than half of US equity volume trades off-exchange. Is that a problem?
For whom?
\end{interviewq}

\begin{interviewq}[$\star\star\star$ \iqroles{researcher, trader}]\label{iq:m1:us-equity-market-structure:6}
The regulator proposes to repeal the \omterm{def:m1:us-equity-market-structure:opr}{order protection rule}. Predict three
consequences for market structure and say how you would position a
market-making firm for each.
\end{interviewq}
